% Fragmento conservado; véase MANIFIESTO_ANTECEDENTES_LECTURA.json.
\section{Generación coinductiva de \texorpdfstring{\(\pi,\varphi,e\)}{pi, phi, e} a profundidad arbitraria}
\label{sec:generacion}

La emisión de seis bloques es una etapa finita de la construcción, no una expansión decimal completa. Su función consiste en producir regiones con multiplicidad, orientación y reglas de transporte. La prolongación debe conservar esos datos y determinar una colección compatible de prefijos de longitud arbitraria. La palabra \emph{generación} designa esta operación; \emph{genealogía} designa la reconstrucción de sus dependencias. Ambas nociones intervienen, pero cumplen funciones distintas.

\subsection{Órbitas, regiones y caracteres de clausura, propagación y autoescala}

Sobre una palabra \(w=(w_1,\ldots,w_6)\), definimos
\[
 r(w)=(w_3,w_1,w_2,w_5,w_6,w_4),\qquad
 s(w)=(w_4,w_5,w_6,w_1,w_2,w_3).
\]
Se cumple \(r^3=s^2=I\) y \(srs=r^{-1}\); por tanto estas permutaciones realizan el grupo diedral \(D_3\). La acción conserva el catálogo y las multiplicidades y conmuta con \(\Psi\). El cociente de las \(243\) palabras visibles tiene \(43\) órbitas: treinta y ocho de cardinal seis y cinco de cardinal tres.

La clase de clausura se reconoce por un predicado de fibra: sus seis orientaciones tienen cinco emisiones por palabra, multiplicidad total \(1008\) y distribución \(432+4\cdot144\). Existe una sola órbita con esas propiedades en el catálogo. Para precisar las otras dos selecciones, escribimos \(f(w)=\#\Psi^{-1}(w)\), \(m(w)=\sum_{U\in\Psi^{-1}(w)}\operatorname{mult}(U)\) y \(\nu(w)=(n_0,n_1,n_2)\). La clase diagonal aditiva de una emisión es \(d(U)=[i+j]_9\) en los índices de cero a ocho; su firma aditiva conserva esta clase y la orientación \(ES\) o \(NO\). Sea \(\mathcal D(\mathcal O)\) el conjunto de esas clases para todas las emisiones sobre una órbita.

Entre las órbitas de tamaño seis, el predicado conjunto
\[
 f(w)=1,\qquad m(w)=144\quad(w\in\mathcal O),\qquad
 \mathcal D(\mathcal O)=\{0,3,6\}
\]
deja exactamente seis órbitas. En este sector, el censo \(\nu=(2,3,1)\), igual al de clausura, selecciona únicamente la propagación; el censo \(\nu=(2,2,2)\) selecciona únicamente la autoescala. El soporte diagonal no puede omitirse: sin él hay cuatro órbitas unipuntuales equilibradas de multiplicidad \(144\), no una. Finalmente, la existencia de una emisión con \(d=6\) y orientación \(ES\) selecciona un representante de cada órbita y conduce a
\begin{equation}
 w^{\rm cl}=010211,\qquad w^{\rm pr}=201101,\qquad w^{\rm au}=121200.
 \label{eq:palabras-regionales}
\end{equation}
Estos símbolos designan primero regiones del catálogo. La identificación escalar se demostrará mediante sus lectores; no se eligen porque coincidan con cifras de una constante conocida.

La palabra de clausura \(010211\) tiene las cinco preimágenes siguientes:
\begin{center}\small
\begin{tabular}{@{}lll r@{}}\toprule
Emisión de seis bloques&Firma multiplicativa&Papel&Multiplicidad\\\midrule
\(501,614,498,169,272,272\)&\(555555\)&transversal&432\\
\(810,923,870,169,272,272\)&\(339555\)&orientación positiva&144\\
\(810,983,810,169,272,272\)&\(393555\)&orientación positiva&144\\
\(870,923,810,169,272,272\)&\(933555\)&orientación positiva&144\\
\(870,923,810,418,674,521\)&\(933717\)&retorno orientado&144\\\bottomrule
\end{tabular}
\end{center}
La multiplicidad \(432\) y la firma constante \(555555\) caracterizan la región transversal. No identifican la posición central \((5,5)\) de APP. Confundir una firma de seis ventanas con una posición del soporte eliminaría la distinción entre región numérica y estructura electrónica.

La microfibra posee además una morfología bipartita precisa. Cada emisión es un producto de una clase de cursor aditivo y una clase multiplicativa. Las tres regiones positivas comparten la misma clase aditiva de dieciocho condiciones; juntas contienen tres clases multiplicativas de ocho condiciones cada una. La región transversal posee un producto \(18\times24\); las tres positivas reunidas, otro \(18\times24\); y el retorno, un producto \(18\times8\). La descomposición por emisiones tiene cinco partes, mientras la descomposición por componentes conexas del grafo bipartito tiene tres. Ambas conservan las mismas \(1008\) semillas sin reducir la forma a un cardinal.

\subsection{Transporte finito y frontera de prolongación}

Los endomorfismos \(L_0,L_1\) del espacio de palabras ternarias actúan antes de la realización arquimediana. Su determinación comienza por las transiciones del estado aritmético, no por la prescripción de dos matrices. Si \(U_t\) es el paso compuesto de selección, transporte y actualización, el bloque observado satisface
\[
 b_j^\chi=\Pi_6(x_j^\chi),\qquad
 x_{j+1}^\chi=U_{9j+9}\cdots U_{9j+1}(x_j^\chi),
 \qquad \chi\in\{\mathrm{cl},\mathrm{pr},\mathrm{au}\}.
\]
Se forman seis transiciones por cada régimen: dos consecutivas en cada una de las tres regiones. Los pares de matrices de orígenes y destinos obtenidos en el registro son
\[
\begin{array}{c|c|c|c}
X_0&Y_0&X_1&Y_1\\\hline
010211&012222&010211&002111\\
012222&010211&002111&110221\\
201101&121221&102011&012222\\
121221&102011&012222&102011\\
121200&112202&121020&010210\\
112202&121020&010210&010200
\end{array}
\]
Cada cadena de seis dígitos representa una fila en \(\FF_3^6\). Puesto que \(\det X_0=1\) y \(\det X_1=-1\) en \(\FF_3\), las ecuaciones \(X_aL_a=Y_a\) determinan unívocamente \(L_a=X_a^{-1}Y_a\). En esa coordenatización de vectores fila resultan
\[
 L_0=\begin{pmatrix}
 2&2&2&1&2&1\\2&1&2&2&1&1\\1&0&0&1&0&2\\
 0&0&1&1&1&0\\2&2&2&0&2&1\\2&1&2&1&0&0
 \end{pmatrix},\quad
 L_1=\begin{pmatrix}
 2&1&2&0&2&2\\1&2&1&1&2&0\\1&2&1&1&1&2\\
 0&1&0&0&2&1\\2&0&2&1&0&1\\0&2&2&2&1&1
 \end{pmatrix}\quad(\bmod3).
\]
El calendario \((L_0,L_0,L_1,L_1)\) produce cuatro nuevos bloques de seis componentes. Sus concatenaciones forman
\[
 w_6\longrightarrow w_{12}\longrightarrow w_{18}
 \longrightarrow w_{24}\longrightarrow w_{30}.
\]
Los subíndices de \(w\) indican longitud acumulada. La frontera \(R_{36}\) conserva la estructura terminal y abre la relación de prolongación; no se renombra como un último bloque periódico que pudiera repetirse indefinidamente.

La identidad \(L=X^{-1}Y\) demuestra unicidad una vez fijadas las transiciones; no sustituye la producción de esas transiciones por \(U_t\). La procedencia de su registro y el examen de los dieciséis calendarios binarios se documentan en el suplemento técnico. El calendario indicado es el único de esos dieciséis que reproduce conjuntamente las tres sucesiones de cinco bloques.

\subsection{Cinco regiones de clausura y compensación angular}

La región de clausura contiene una componente transversal y cuatro acompañantes. El lector simétrico actúa en el espacio de sus cinco posiciones mediante
\[
 P_5=\frac15\mathbf1\mathbf1^{\mathsf T},\qquad
 \mathbf1=(1,1,1,1,1)^{\mathsf T}.
\]
La invariancia \(P_5\lambda=\lambda\) y la conservación \(\mathbf1^{\mathsf T}\lambda=1\) fuerzan \(\lambda=\mathbf1/5\). Esta normalización distribuye la unidad entre posiciones del lector; no sustituye las multiplicidades de semillas \(432,144,144,144,144\) por frecuencias uniformes. La coordenatización elíptica del lector realiza cada coordenada mediante \(1+\lambda_jJ\); su pendiente es \(\lambda_j\), por lo que la pendiente primitiva resulta \(1/5\). Se hace explícita así la regla que pasa de una coordenada normalizada a una pendiente, en lugar de identificar ambas nociones sin mapa. El lector compone las cuatro posiciones acompañantes por la ley
\begin{equation}
 a\oplus b=\frac{a+b}{1-ab}.
 \label{eq:suma-pendientes}
\end{equation}
Esta ley se obtiene al multiplicar \((1+aJ)(1+bJ)\) en el régimen \(J^2=-I\) y dividir por su componente escalar. Por ello la composición de pendientes es una realización de la estructura cuadrática de TRIT, no una suma ordinaria de cuatro fracciones.

\begin{proposition}[Compensador de la clausura]
La composición de cuatro pendientes \(1/5\) es \(120/119\). La condición de retorno a pendiente uno determina un único compensador positivo de la forma \(1/q\), con \(q>1\), y da \(q=239\).
\end{proposition}
\begin{proof}
La primera duplicación da \((1/5)\oplus(1/5)=5/12\); la segunda, \((5/12)\oplus(5/12)=120/119\). La condición
\[
 \frac{120/119-1/q}{1+(120/119)/q}=1
\]
equivale a \((120-119)q=120+119\), por lo que \(q=239\).
\end{proof}

El carácter de clausura resultante tiene realización
\begin{equation}
 \Pi_{\rm cl}=4\left(4\arctan\frac15-\arctan\frac1{239}\right).
 \label{eq:lector-pi}
\end{equation}
La identidad angular de Machin reconoce \(\Pi_{\rm cl}=\pi\). El sentido de la construcción es preciso: la pendiente y el número de composiciones están declarados por el lector de la pentafibra; la ecuación de compensación fuerza \(239\). Una vez producido este carácter, su evaluación por series alternadas no selecciona de nuevo la órbita de semillas.

Para \(0<x<1\), las sumas alternadas
\[
 A_n(x)=\sum_{j=0}^n\frac{(-1)^jx^{2j+1}}{2j+1}
\]
encierran \(\arctan x\) entre dos truncamientos consecutivos, con diferencia \(x^{2n+3}/(2n+3)\). Aplicadas a las dos pendientes de \eqref{eq:lector-pi}, proporcionan extremos racionales de error explícito. El valor numérico se calcula así como consecuencia del carácter exacto, sin introducir un prefijo objetivo.

\subsection{Lector de propagación y normalización de la unidad}

La región de propagación se realiza mediante una colección formal \(P_t\in\mathcal A[[t]]\), donde \(\mathcal A\) es un álgebra conmutativa unitaria sobre \(\QQ\). Su composición satisface \(P_{t+s}=P_tP_s\), con unidad \(P_0=1\), y el transporte elemental orientado fija el generador escalar \(P'_0=+1\). La condición fija su signo, no sólo su magnitud. Escribiendo
\[
 P_t=\sum_{n\ge0}a_nt^n,
\]
la normalización determina \(a_0=a_1=1\). La comparación del coeficiente de \(s^nt\) en \(P_{s+t}=P_sP_t\) da \((n+1)a_{n+1}=a_na_1\); equivalentemente, \(P'_t=P_t\). Así,
\begin{equation}
 (n+1)a_{n+1}=a_n,\qquad a_n=\frac1{n!}.
 \label{eq:propagacion}
\end{equation}
El carácter en una unidad de parámetro es \(\Pi_{\rm pr}=P_1\), reconocido posteriormente como \(e\).

La normalización del generador importa. La ley semigrupal sin ese dato admitiría \(P_t=e^{ct}\) para distintos \(c\); por ello no se omite la condición que fija la unidad del parámetro en el lector. Esta condición es anterior a la comparación decimal y pertenece a la exposición de la regla operacional, no a la elección de un valor experimental.

\begin{proposition}[Cotas racionales de propagación]
Para \(n\ge1\), sea \(E_n=\sum_{j=0}^n1/j!\). Entonces
\[
 E_n<\Pi_{\rm pr}<E_n+\frac{n+2}{(n+1)(n+1)!}.
\]
\end{proposition}
\begin{proof}
El primer término omitido es \(1/(n+1)!\). Cada cociente posterior es a lo sumo \(1/(n+2)\). La suma geométrica mayorante de la cola es
\[
 \frac1{(n+1)!}\frac1{1-1/(n+2)}
 =\frac{n+2}{(n+1)(n+1)!}.
\]
La desigualdad estricta se sigue de la disminución de los cocientes sucesivos.
\end{proof}

\subsection{Incidencia modal y generación de la autoescala}

El selector trítico produce el ciclo \((+1,-1,0)\). En los modos \(\Sigma,\Pi\), la regla es \(+1:m\mapsto\Sigma\), \(-1:m\mapsto\Pi\), \(0:m\mapsto m\). Con cierre cíclico, el modo observado es \((\Sigma,\Pi,\Pi)\). Sus aristas son \(\Sigma\to\Pi\), \(\Pi\to\Pi\) y \(\Pi\to\Sigma\). Por tanto su incidencia, en ese orden de modos, es
\begin{equation}
 F_{\rm av}=\begin{pmatrix}0&1\\1&1\end{pmatrix}.
 \label{eq:fav}
\end{equation}
Las cuatro entradas se obtienen del calendario, no de una ecuación que haya recibido previamente el valor de \(\varphi\). La repetición del calendario da conteos \(3F_{\rm av},9F_{\rm av},36F_{\rm av}\) en longitudes nueve, veintisiete y ciento ocho. Estos conteos no son las potencias de la matriz.

\begin{theorem}[Carácter positivo de autoescala]
La incidencia \eqref{eq:fav} posee un único rayo propio estrictamente positivo. Al normalizarlo como \((1,x)^{\mathsf T}\), su coordenada \(x\) satisface \(x^2-x-1=0\), \(x>0\), y es \(\varphi=(1+\sqrt5)/2\).
\end{theorem}
\begin{proof}
La ecuación propia da \(x=\lambda\) y \(1+x=\lambda x\). Eliminar \(\lambda\) produce el polinomio. De sus dos raíces sólo una es positiva. Como \(F_{\rm av}^2\) tiene todas sus entradas positivas, el rayo positivo es único y simple. La expresión radical reconoce la coordenada obtenida; no participa en la producción de \(F_{\rm av}\).
\end{proof}

Con \(F_0=0,F_1=1,F_{n+1}=F_n+F_{n-1}\), la incidencia induce, para \(n\geq1\),
\[
 F_{\rm av}^n=\begin{pmatrix}F_{n-1}&F_n\\F_n&F_{n+1}\end{pmatrix}.
\]
Los cocientes consecutivos \(F_{n+1}/F_n\) alternan alrededor de \(\varphi\). La identidad de determinante
\[
 F_{n+1}^2-F_nF_{n+2}=(-1)^n
\]
da una anchura exacta \(1/(F_nF_{n+1})\) entre dos cocientes consecutivos. Su crecimiento hace arbitrariamente pequeña esta horquilla racional.

\subsection{Cilindros compatibles y profundidad arbitraria}

La realización de un bloque ternario \(b=(b_1,\ldots,b_6)\) es el entero
\[
 \nu(b)=\sum_{j=1}^6b_j3^{6-j}\in\{0,\ldots,728\}.
\]
Una sucesión de bloques define prefijos enteros por \(N_{n+1}=729N_n+\nu(b_{n+1})\), comenzando en \(N_0=m\), donde \(m\) es la parte entera determinada por el lector. Las horquillas anteriores sitúan los caracteres de clausura, propagación y autoescala en \((3,4)\), \((2,3)\) y \((1,2)\), respectivamente; por tanto, sus valores iniciales son \(m=3,2,1\). Los bloques siguientes refinan la parte fraccionaria. El cilindro correspondiente es
\begin{equation}
 I_n=\left[\frac{N_n}{729^n},\frac{N_n+1}{729^n}\right],
 \qquad |I_n|=729^{-n}.
 \label{eq:cilindro}
\end{equation}
La convención semiabierta se utiliza para seleccionar un representante cuando un valor cae en una frontera. Para el paso al límite se utilizan las clausuras de estos intervalos y se identifican las dos expansiones fronterizas equivalentes. Esta precisión impide suponer que toda intersección de intervalos semiabiertos encajados sea automáticamente no vacía.

\begin{theorem}[Determinación de la realización escalar]
Una historia de bloques compatible determina un único límite de los extremos izquierdos de \eqref{eq:cilindro}. Si un carácter exacto dispone de horquillas racionales de anchura arbitrariamente pequeña y se separa de la frontera de cada cilindro decimal examinado, cada prefijo decimal de longitud finita se determina después de un cálculo finito. Los prefijos así obtenidos son compatibles entre sí.
\end{theorem}
\begin{proof}
La extensión de un bloque conserva el prefijo anterior por división euclídea. Los cilindros cerrados están encajados y su diámetro tiende a cero; sus extremos izquierdos son de Cauchy y tienen un único límite. Para una longitud decimal fijada, un valor separado de las fronteras de la partición posee distancia positiva a ellas. Una horquilla suficientemente estrecha queda contenida en una única celda y determina su prefijo. La unicidad de la celda y el encajamiento de las particiones prueban la compatibilidad por truncamiento. En los puntos con expansión terminal debe añadirse la decisión exacta de frontera y su convención de representación; no se deduce esa decisión de una anchura numérica solamente.
\end{proof}

\begin{corollary}[Obtención finita y compatibilidad de los prefijos regionales]
\label{gen:prefijos-regionales}
Para cada una de las realizaciones \(\pi,\varphi,e\) y para toda longitud
decimal finita, las horquillas construidas determinan el prefijo correspondiente
mediante un cálculo finito. Los prefijos obtenidos a distintas profundidades
son compatibles por truncamiento.
\end{corollary}
\begin{proof}
Los caracteres de clausura, propagación y autoescala poseen las horquillas
anteriores. Tras su reconocimiento, la irracionalidad de \(\pi,e,\varphi\)
garantiza su separación de las fronteras racionales de cada partición decimal
finita. Se aplica el teorema precedente a cada carácter.
\end{proof}
La profundidad arbitraria expresa el cuantificador del corolario: para todo
horizonte finito existe un cálculo que determina su prefijo. El objeto
matemático es la colección compatible completa; cada ejecución obtiene una
proyección finita de ella.

\subsection{Involución especular y sincronización de las regiones}

La conexión nonádica transporta conjuntamente los tres caracteres. En la coordenatización de la novena fase, el eje de autoescala y la suma ternaria de clausura y propagación permanecen fijados bajo el intercambio especular. La inversión helicoidal cambia el sentido del transporte; el espejo de los dos canales cambia su orientación relativa. Son involuciones diferentes, aunque ambas actúen sobre el mismo sistema de trayectorias.

Las coincidencias de censo ilustran esa coordinación:
\begin{center}
\begin{tabular}{@{}c rrr l@{}}\toprule
Fase&\(N_\pi\)&\(N_e\)&\(N_\varphi\)&Igualdad de conteos\\\midrule
4&207&207&122&\(N_\pi=N_e\)\\
5&39&151&151&\(N_e=N_\varphi\)\\
6&110&110&69&\(N_\pi=N_e\)\\
7&81&29&81&\(N_\varphi=N_\pi\)\\
8&59&59&59&sincronización triple\\
9&43&43&43&sincronización triple\\\bottomrule
\end{tabular}
\end{center}
Estos enteros cuentan coordenadas compatibles de las regiones. No afirman una igualdad entre \(\pi\), \(e\) y \(\varphi\) como números reales. Si
\[
 \partial(a,b,c)=(a-b,b-c,c-a),
\]
su imagen pertenece al retículo \(A_2=\{(u,v,w)\in\ZZ^3:u+v+w=0\}\). En las fases ocho y nueve la componente relativa se anula, pero la componente diagonal cambia de \((59,59,59)\) a \((43,43,43)\). La sincronización relativa coexiste con un desplazamiento de \(-16\) por canal. El cierre que conducirá a \(\alpha\) debe conservar ambas informaciones: coincidencia relativa y evolución de la memoria.

